% 动态执行与反馈多智能体架构图 — 独立编译：xelatex multiagent_architecture.tex
\documentclass[tikz,border=12pt]{standalone}
\usepackage{ctex}
\usetikzlibrary{arrows.meta,calc,positioning}

\begin{document}
\begin{tikzpicture}[
    font=\small,
    hub/.style={
        rectangle,
        rounded corners=6pt,
        draw=blue!70,
        very thick,
        fill=blue!12,
        minimum width=3.65cm,
        minimum height=1.2cm,
        align=center,
        font=\bfseries
    },
    sub/.style={
        rectangle,
        rounded corners=6pt,
        draw=black!55,
        very thick,
        minimum width=3.05cm,
        minimum height=0.95cm,
        align=center,
        font=\bfseries\small
    },
    tool/.style={
        rectangle,
        rounded corners=3pt,
        draw=black!45,
        thin,
        fill=cyan!8,
        minimum height=0.5cm,
        inner xsep=5pt,
        align=center,
        font=\scriptsize
    },
    plat/.style={
        rectangle,
        rounded corners=5pt,
        draw=purple!58,
        thick,
        fill=purple!7,
        minimum width=7.2cm,
        minimum height=0.95cm,
        align=center,
        font=\small\bfseries
    },
    lbl/.style={font=\scriptsize, inner sep=1pt, fill=white, fill opacity=0.93, text opacity=1},
    invoke/.style={{Stealth[length=2.5mm]}-, thick, dashed, draw=orange!78!black},
    solid/.style={{Stealth[length=2.5mm]}-, thick, draw=black!58},
    loop/.style={{Stealth[length=2.5mm]}-, semithick, dashed, draw=red!65!black}
]

% 主控
\node[hub] (runner) at (0, 0) {Emulator Runner Agent\\{\scriptsize\mdseries LangGraph 主控}};

% 主控直连工具
\node[tool, above=0.5cm of runner.west, anchor=south west, xshift=-0.05cm] (t1) {EmulateProject};
\node[tool, above=0.5cm of runner, anchor=south] (t2) {GetMMIOFunctionEmulateInfo};
\node[tool, above=0.5cm of runner.east, anchor=south east, xshift=0.05cm] (t3) {GetFunctionCallsEmulateInfo};
\draw[solid] (t1.south) -- (t1.south |- runner.north);
\draw[solid] (t2.south) -- (runner.north);
\draw[solid] (t3.south) -- (t3.south |- runner.north);

% 子智能体
\node[sub, fill=orange!15, draw=orange!72!black] (fixer) at (-4.75, 0) {Fixer Agent};
\node[sub, fill=green!12, draw=green!58!black] (builder) at (4.75, 0) {Builder Agent};

\node[tool, below=0.4cm of fixer.south, anchor=north, minimum width=3.35cm] (fu) {{\scriptsize\texttt{update\_function\_replacement}}};
\node[tool, above=0.38cm of builder.north, anchor=south] (br) {\texttt{replace\_funcs}};
\node[tool, right=0.28cm of br.east, anchor=west] (bp) {\texttt{build\_project}};
\draw[solid] (fixer) -- (fu);
\draw[solid] (builder) -- (br);
\draw[solid] (br) -- (bp);

% 工具式调用
\draw[invoke] (runner.west) -- node[lbl, above=1pt, sloped] {Fixer 工具} (fixer.east);
\draw[invoke] (runner.east) -- node[lbl, above=1pt, sloped] {Builder 工具} (builder.west);

% 仿真后端
\node[plat] (emu) at (0, -2.75) {MCP $\rightarrow$ Unicorn 固件仿真};
\draw[solid] (runner.south) -- node[lbl, right=1pt] {下发仿真} (emu.north);
\draw[solid] (emu.140) to[bend left=18] node[lbl, pos=0.45, sloped] {日志 / 崩溃 / 结果} (runner.220);

% Fixer 写入与 Builder 读取：经数据管理器（概念带）
\coordinate (dm) at (0, -1.35);
\node[lbl, draw=gray!55, rounded corners=2pt, inner sep=3pt, fill=gray!6] at (dm) {替换数据管理};
\draw[solid, draw=gray!50] (fu.south) |- ($(dm)+(0,-0.12)$);
\draw[solid, draw=gray!50] ($(dm)+(0,-0.12)$) -| (br.south);

% 构建输出回归主控
\draw[loop] (bp.east) -- ++(0.45, 0) |- node[lbl, pos=0.23, right=1pt] {BuildOutput} (runner.north east);

\end{tikzpicture}
\end{document}
